\section{The exact characteristic-\(\ell\) curve criterion for the cubic model}
\label{sec:cubic-characteristic-curve-criterion}

Let \(O\), \(A\), \(J\), \(S\) and the actual normalized cubic
representation \(r_S\) be those of AC-G51--AC-G58.
In particular \(O\) is a complete discrete valuation ring with
finite residue field \(k\) of characteristic
\(\ell\ne2,3,269\).  All retained arithmetic and marked-point
hypotheses of that construction remain in force.
Put
\[
 B=\operatorname{im}(A_0\longrightarrow A/J)
       \ \subset S .
                                                               \tag{G59.1}
\]
The inclusion is the homogeneous contraction proved in AC-G56.
The exact source condition \texttt{M2} says that \(B\) is a
finite \(O\)-module.

A \emph{characteristic-\(\ell\) curve of the actual ring} means
a local \(O\)-algebra map
\[
 f:S\longrightarrow V,\qquad
 V\text{ a complete discrete valuation ring with finite
 residue field},\qquad \varpi V=0 .
                                                               \tag{G59.2}
\]
Its residue field may be a finite extension of \(k\).
No lifting of this curve to characteristic zero is assumed.
Write \(r_f=f_*r_S\), let \(H=\ker\eta=G_K\), and write the
three residual characters on \(H\) as
\[
 \alpha=\overline\chi^{\,j}\overline\gamma,\qquad
 1,\qquad
 \delta=\overline\chi^{\,j}\overline\gamma^{-1},
                    \qquad j=1,2 .
                                                               \tag{G59.3}
\]
The normalized seed has matrix \(I\) at complex conjugation,
and so does every \(r_f\).  The marked difference
\(\overline\chi(g_1)^j
(\overline\gamma(g_1)-\overline\gamma(g_1)^{-1})\)
is nonzero.

\begin{theorem}[AC-G59: an integral curve criterion for the actual \texttt{M2}]
\label{AC-G59}
The following conditions are equivalent.
\begin{enumerate}
\item \(B\) is finite over \(O\), that is, the actual
      \texttt{M2} holds.
\item For every curve (G59.2), the representation \(r_f\)
      is upper triangular or lower triangular on \(H\),
      and its diagonal is the prescribed seed diagonal.
\item In the weighted quotient presentation of \(S\)
      induced by the actual \(W_A\) of AC-G51, every
      curve sends all weight-zero coordinates to zero
      and sends either all positive-weight coordinates
      or all negative-weight coordinates to zero.
\end{enumerate}
The statements concern all actual characteristic-\(\ell\)
curves, including those which detect vertical components.
\end{theorem}
\begin{proof}
\proofstep{The presentation and the zero part.}
Pass the weighted presentation
\(P=O[[z_1,\ldots,z_d]]\to R\) of AC-G56 through
\(R\to S\).
Its kernel is homogeneous: the old kernel is homogeneous,
and the additional ideal is generated by the homogeneous
global and normalization equations.
Call the resulting weighted presentation \(W_S\).
Its formal degree-zero part is exactly \(B\), since
both are the image of the homogeneous series \(P_0\).
The ring \(B\) is complete Noetherian local with residue
field \(k\), by the proof of AC-G58-algebra applied to
this quotient presentation.

Every coordinate of \(W_S\) is an embedded universal
entry difference, including the restricted local
coordinates; this is the exact construction in
(G51.7)--(G51.8).
As proved in AC-G51, the embedded entries are actual
entries of \(r_S(g)\), entries of \(r_S(cg)\), or zero,
with the same weight.
Every coordinate belongs to the maximal ideal of \(S\).

\proofstep{Integral finiteness kills the zero-part deviations.}
Assume condition 1 and fix \(f\).
The map \(B\to V\) factors through the finite-dimensional
local \(k\)-algebra \(B/\varpi B\).
Its maximal ideal is nilpotent.  Since \(V\) is a domain,
that ideal maps to zero.
Thus every \(b\in B\) takes its residue-field value,
and every element of \(\mathfrak m_B\) maps to zero.
In particular all weight-zero coordinates map to zero.

Suppose a positive coordinate \(z_i\), of weight \(a>0\),
and a negative coordinate \(z_j\), of weight \(-b<0\),
both have nonzero images.
The monomial \(z_i^b z_j^a\) has weight zero and lies in
\(\mathfrak m_B\).
Its image would be both zero and a product of nonzero
elements of the domain \(V\), a contradiction.
Consequently one entire sign of the coordinates is
annihilated.  This proves condition 3.

\proofstep{Coordinate vanishing and triangularity.}
Condition 2 implies condition 3 by the entry description
of the coordinates.
Conversely, a local homomorphism \(S\to V\) is continuous
for the maximal-ideal topologies: it sends
\(\mathfrak m_S^N\) into \(\mathfrak m_V^N\).
If all positive coordinates vanish, every monomial of
positive weight vanishes, and so does every formal
series of positive weight, by convergence in \(V\).
The same argument applies with signs reversed.
Every nonconstant zero-weight monomial contains either
a weight-zero coordinate or coordinates of both signs.
Hence its image is zero as well.
The homogeneous universal entries therefore have the
required triangular vanishing, and the diagonal entry
differences, which have weight zero and zero constant
term, vanish.
This proves condition 2.
For elements outside \(H\), right-multiplying by \(c\) changes
neither the matrix of \(r_f\) nor that of the seed,
because their matrices at \(c\) are \(I\).

\proofstep{The characteristic-\(\ell\) curve test includes vertical torsion.}
The precise source theorem
\path{CurveDVRSpecialization/FinCrit.lean},
\texttt{UpperTheory.V.Pr.finCrit}, applies to \(W_S\).
Its quantified hypotheses are: \(O\) complete discrete
valuation with finite residue field; \(S\) an
\texttt{IsCO} \(O\)-algebra; all presentation coordinates
in \(\mathfrak m_S\); and, for every local map to an
\texttt{IsCD} \(O\)-algebra \(V\) annihilating
\(\mathfrak m_O\), precisely the two coordinate
vanishing assertions in condition 3.
Here \texttt{IsCD} is exactly the complete discrete
valuation ring with finite residue field in (G59.2);
its \(O\)-algebra map is local.
All these hypotheses have just been verified.
The theorem concludes that the formal degree-zero part
of \(W_S\), hence \(B\), is finite over \(O\).

For clarity about the integral content of this applied
theorem, its proof does not merely inspect a generic
fiber.  The proof uses \texttt{NicePrimeEx.npe\_dim\_one}
and its own \texttt{fc\_point} to detect a nonnilpotent
zero-weight monomial modulo \(\varpi\) on a
one-dimensional quotient and then on a complete
discrete valuation ring.
The two vanishing assertions exclude that curve.
It follows that a power of \(\mathfrak m_B\) lies in
\(\varpi B\), so \(B/\varpi B\) is finite; complete
module Nakayama gives the asserted integral finiteness.
This is the source's integral argument, with its
literal curve hypothesis, applied to the cubic model.
\end{proof}

\begin{lemma}[AC-G59-polynomials: constant characteristic polynomials for a cubic curve]
\label{AC-G59-polynomials}
Fix a curve (G59.2).  Suppose that for every \(h\in H\)
all coefficients of the characteristic polynomial of
\(r_f(h)\) are algebraic over \(k\).
Then condition 2 of AC-G59 holds for this curve.
In fact
\[
 \det(T-r_f(h))
       =(T-\alpha(h))(T-1)(T-\delta(h))
                         \qquad(h\in H)
                                                               \tag{G59.4}
\]
as an equality in \(V[T]\).
Consequently \texttt{M2} is also equivalent to this
algebraicity condition on every actual curve.
\end{lemma}
\begin{proof}
\proofstep{Algebraic values with fixed residue are constant.}
If \(x\in V\) is algebraic over \(k\) and reduces to
\(a\in k\), then \(x=a\).
Indeed \(k[x]\) is a finite field extension of \(k\),
inside the domain \(V\).
If \(x-a\ne0\), its inverse belongs to that field
and therefore to \(V\), contradicting
\(x-a\in\mathfrak m_V\).
Apply this to every characteristic coefficient:
the actual residual representation prescribes its
residue in \(k\).
This proves (G59.4).

\proofstep{A split character polynomial cannot be irreducible.}
We record the elementary representation-theoretic
argument used below.
Over an algebraically closed field, an irreducible
representation \(W\) of dimension \(n>1\) cannot have
the same characteristic polynomials as a direct sum
of \(n\) characters.
Otherwise their traces would agree on every group
element.
The image algebra of the group algebra on
\(W\oplus\bigoplus_i k(\tau_i)\) is finite dimensional
and semisimple: the module is a faithful direct sum
of simple modules, so its radical acts as zero and
is zero.
Its simple matrix factor for \(W\) is distinct from
every one-dimensional character factor.
In that product of matrix algebras take an element
acting as \(E_{11}\) on \(W\) and as zero on every
character summand.
It is a finite linear combination of group operators.
The asserted trace equality, extended linearly,
would say \(1=0\).
This proves the assertion in every characteristic.
It is the same image-algebra argument as AC-G38.
The pinned source's seed-independent
\path{CurveDVRSpecialization/NotIrrOfChar.lean},
\texttt{notIrrOfChar}, supplies the corresponding
general characteristic-polynomial lemma.

\proofstep{Use the actual normalized shape.}
Extend scalars to an algebraic closure \(L\) of
\(\operatorname{Frac}(V)\).
The exact theorem
\path{UnitaryThreeGroups/ExistsCoordOfNotMatIrr.lean},
\texttt{shape\_cases}, depends on the semidirect
multiplier, \(2\ne0\), normalization at \(c,g_1\),
and the nonzero marked \texttt{zeta}; it does not
require that the residual tuple be an adjoint tuple.
All its hypotheses hold here.
It gives Type I, upper triangular, lower triangular,
or Type II.

Type I is excluded by (G59.4) and the preceding
image-algebra argument in dimension three.
In the Type II case the endpoint block is an
irreducible two-dimensional representation of \(H\),
and the middle entry is a character.
For each \(h\), that middle entry is a root of the
polynomial (G59.4), hence is algebraic over \(k\).
It belongs to \(V\) and reduces to \(1\), so it equals
\(1\) by the first step.
Canceling \(T-1\) from the block characteristic
polynomial gives
\[
       \det(T-r_f(h)|_{\{0,2\}})
                       =(T-\alpha(h))(T-\delta(h)).
                                                               \tag{G59.5}
\]
The same argument in dimension two excludes the
irreducibility of that block.
Thus the representation is upper or lower triangular
over \(L\), hence over \(V\) by injectivity.

Each diagonal entry of a triangular \(r_f(h)\) is
a root of (G59.4).
It is algebraic over \(k\) and has its prescribed
residual coordinate value, so equals that value.
This gives precisely the seed diagonal.
Multiplication by \(c\) gives the same assertion
outside \(H\).
The converse needed for the final equivalence is
immediate from triangularity with this diagonal,
and AC-G59 completes the proof.
\end{proof}

\begin{corollary}[AC-G59-finite-index: only the actual curves need the arithmetic input]
\label{AC-G59-finite-index}
It suffices for \texttt{M2} that for every actual curve
\(f:S\to V\) there be a finite-index subgroup
\(U_f\) of the absolute Galois group such that
\[
 \det(T-r_f(h))=\det(T-\overline r_*(h))
                         \qquad(h\in U_f\cap H).
                                                               \tag{G59.6}
\]
The subgroup may depend on the curve.
In fact the resulting conclusion is (G59.4) on all
of \(H\), for every curve.
This is a sufficient test on the actual selected
deformation ring, without an assertion about all
ordinary representations outside its support.
\end{corollary}
\begin{proof}
Fix \(f\) and \(h\in H\).
Finite index supplies \(n>0\) with \(h^n\in U_f\).
On \(H\) the matrix component is a linear
representation, so \(r_f(h^n)=r_f(h)^n\).
By (G59.6) all eigenvalues of this power are in \(k\).
Every eigenvalue of \(r_f(h)\) therefore satisfies
an equation \(X^n-a=0\) with \(a\in k\).
It is algebraic over \(k\), and so are the elementary
symmetric functions giving the characteristic
coefficients.
This argument does not require \(n\) prime to
\(\ell\).
AC-G59-polynomials proves the full characteristic
polynomial identity and triangularity for this
curve; AC-G59 then gives \texttt{M2}.
\end{proof}

\paragraph{The exact source supply still to be transported.}
The definition \texttt{FinEnd} in
\path{CurveDVRSpecialization/RatTorsionFree.lean}
uses one finite-index subgroup for all complete-DVR
ordinary representations annihilating \(\mathfrak m_O\).
The full proof
\path{CurveDVRSpecialization/GfResZpow.lean},
\texttt{finA\_of\_end}, passes this hypothesis through
\texttt{constOfChar} and the integral
\texttt{finCrit}, yielding the finite
\texttt{weightZeroSubalgebra} used in
\texttt{modelIdeal\_finiteness\_of\_weightZero\_finite}.
The latter then supplies the literal \texttt{M2}.
The present proof makes the cubic tuple substitution
in the constant-polynomial argument explicit and
restricts the required test to curves of the actual
ring.  It does not assert that this arithmetic test
has been proved.

The supplied arithmetic bridge
\path{CurveDVRSpecialization/FinEndOfFinPt.lean},
\texttt{finEnd\_of\_finPt}, still requires a
\texttt{GalLink} and \texttt{FinPt} for the prescribed
source frame.  Its proof constructs an
\texttt{exists\_nrep3} and uses the explicit identity
between the frame characters and
\(\overline\Psi^{\,\mathrm{wt1}(i)}\).
The assembly in
\path{ModularityArgument/ThresholdArguments/DenseThresholdCases/GdEndOf.lean},
\texttt{gdEnd\_of\_frame\_linkedEnd\_downCsdK\_descent},
retains this same link and the rational Dirichlet datum.
An actual cubic curve version of that arithmetic
input remains to be supplied.
The integral ordinary flags of AC-G57-integral
address its local ordinary condition; they do not
prove constancy of the global curve representation.

\paragraph{Why generic-fiber finiteness is insufficient.}
The complete local ring
\[
                 O[[u]]/(\varpi u)
                                                               \tag{G59.7}
\]
has generic fiber \(\operatorname{Frac}(O)\), but
its reduction modulo \(\varpi\) is \(k[[u]]\).
It is therefore not finite over \(O\).
The local curve to \(k[[t]]\), given by
\(\varpi\mapsto0,\ u\mapsto t\), detects the missing
vertical part.
This is a counterexample to the generic-fiber
inference only, not a claimed countermodel of the
full cubic deformation problem.

\paragraph{Geometric dictionary.}
The curve criterion tests the special fiber of the
formal deformation space by complete discrete
valuation rings.
Its finite local systems and their restrictions
retain the already audited SGA~1, Expos\'e~V,
\S\S4 and~7 interpretation.
The integral criterion, constant-polynomial
argument and selected-ring scope are the explicit
proofs above; the finite-cover dictionary itself
does not supply the remaining arithmetic
\texttt{FinPt} or \texttt{FinEnd} input.
